\section{Problem Formulation}
\label{sec:problem-formulation}

\paragraph{Search space.} Let $[n] \triangleq \{1, \dots, n\}$. A genotype
$x = (x_1, \dots, x_n) \in \Lambda = \Lambda_1 \times \cdots \times \Lambda_n$ assigns one categorical value to
each coordinate. The first coordinates choose one operation per edge of the cell graph (e.g.\
\textsc{conv}$3{\times}3$, \textsc{skip}, \textsc{none}); the remaining coordinates choose training
hyperparameters such as learning rate or weight decay. Continuous hyperparameters are discretised into a fixed
grid before search (Section~\ref{sec:proposed-optimizer}), so $\Lambda$ is entirely categorical.

\paragraph{Validity.} A decoding function $D \colon \Lambda \to \mathcal{N}$ maps a genotype to a network. Not
every genotype decodes to a valid cell (e.g.\ no path from input to output). We write the constraint as
$g(x) \leq 0$ and the feasible set as $\Lambda^* = \{x \in \Lambda : g(x) \leq 0\}$.

\paragraph{Objectives.} Each $x \in \Lambda^*$ is assessed by
\begin{align}
    f_1(x) &= \text{validation error of } D(x) \quad \text{(minimised)}, \\
    f_2(x) &= \text{computational cost of } D(x) \quad \text{(minimised)}.
\end{align}
The cost measure is the one each benchmark exposes: model size in megabytes on JAHS-Bench-201, and tabulated
total training-and-validation time on NAS-HPO-Bench-II, which is read from the same lookup-table row as $f_1$
(one fixed recorded seed) and is therefore not cheaper to obtain. Both objectives are always queried at the highest fidelity $r_K$ the
benchmark tabulates. We seek the Pareto front
\[
    \mathcal{P}^* = \bigl\{ x \in \Lambda^* : \nexists\, x' \in \Lambda^*,\;
    \mathbf{f}(x') \preceq \mathbf{f}(x) \wedge \mathbf{f}(x') \neq \mathbf{f}(x) \bigr\},
\]
where $\mathbf{f} = (f_1, f_2)$ and $\preceq$ is component-wise $\leq$.

\paragraph{Budget.} The cost of search is the number of full evaluations of $\mathbf{f}$. Both benchmarks answer
deterministically, so each configuration is evaluated once; randomness across the $R$ independent runs of a
method comes from initialisation and stochastic operators, not from repeated queries.

\paragraph{Relative surrogate.} After $t$ full evaluations we hold
$\mathcal{H}_t = \{ (x^{(i)}, \mathbf{f}(x^{(i)})) \}_{i=1}^{t}$. Rather than regressing $f_1$ from an encoded
genotype, as NSGANetV2 does~\citep{lu2020nsganetv2}, P3Net's surrogate is \emph{relative} and
\emph{linkage-aware}. For a parent $x \in \Lambda^*$ and a candidate $x'$ that differs from $x$ only on a
linkage subset $F \subseteq [n]$ of the current dependency model ($x'_i = x_i$ for all $i \notin F$), it
predicts the effect of that local change,
\[
    \hat{\delta}_F(x, x') \approx f_1(x) - f_1(x'),
\]
trained on pairwise differences derived from $\mathcal{H}_t$. A positive $\hat\delta_F$ predicts an improvement.
